% ResearchBench — 30-minute talk draft
% Build: pdflatex researchbench_talk.tex  (twice for TOC)
% ~22 content slides ≈ 30 min at ~80s/slide. Cut slides marked [OPTIONAL] if running long.
\documentclass[aspectratio=169]{beamer}
\usetheme{Madrid}
\usecolortheme{seahorse}
\usepackage{booktabs}
\usepackage{graphicx}
\setbeamertemplate{navigation symbols}{}
\setbeamerfont{frametitle}{series=\bfseries}

\title{ResearchBench: Paper Replication as the Benchmark\\for Autonomous AI Research Agents}
\subtitle{Rubric-scored replication of 20 inference-optimization papers on a single H100}
\author{Annika Kaul Singh \\ \texttt{research@infinity.inc}}
\date{July 9, 2026}

\begin{document}

\begin{frame}
  \titlepage
\end{frame}

\begin{frame}{Outline}
  \tableofcontents[hideallsubsections]
\end{frame}

% =====================================================================
\section{Why paper replication}
% =====================================================================

\begin{frame}{Paper replication is a superb eval target for research agents}
  \begin{itemize}
    \item \textbf{Novel research is mostly re-composition.} The sub-tasks of new work
          --- environment bring-up, data prep, method implementation, benchmarking,
          analysis --- replicate the vast majority of sub-tasks in pre-existing research.
    \item An agent that can \emph{faithfully replicate} a paper has demonstrated the
          skill set needed to \emph{execute} novel research.
    \item Ground truth exists: the paper's own reported numbers.
    \item Difficulty is tunable: from structural scaffolding to exact numeric
          reproduction under tolerance.
  \end{itemize}
  \vfill
  \begin{block}{The question this benchmark answers}
    \emph{Which coding agent --- and which harness --- should you trust to do
    autonomous AI research?}
  \end{block}
\end{frame}

\begin{frame}{From binary replication to a score: rubrics}
  \begin{columns}[T]
    \column{0.52\textwidth}
      \textbf{2024 --- PaperBench:} explicit task decompositions (``rubrics'')
      turn replication from a binary judgment into graded sub-tasks.
      Rubrics were \emph{hand-authored} by experts.
      \vspace{0.5em}

      \textbf{2026 --- what changed:} terminal agents now \emph{auto-generate}
      both the decomposition and the per-task correctness tests.
      \begin{itemize}
        \item rubric generator + critic (two-pass, convergence loop)
        \item verification contracts with numeric tolerances
        \item shared structural/behavioral test suites
      \end{itemize}
    \column{0.44\textwidth}
      \begin{block}{Replication as a score}
        Each paper $\rightarrow$ a tree of graded leaves.
        \[
          \text{score} = \frac{\#\,\text{leaves verified pass}}{\#\,\text{achievable leaves}}
        \]
        A leaf passes only when an \emph{executable verifier} attributes it ---
        not when the agent claims it.
      \end{block}
  \end{columns}
\end{frame}

\begin{frame}{Pipeline: fully automated rubric generation and grading}
  \centering
  \small
  \texttt{arxiv URL $\rightarrow$ ingest $\rightarrow$ triage $\rightarrow$ rubric generator $\rightarrow$ rubric critic}\\[2pt]
  \texttt{$\rightarrow$ feasibility / unblockers / rubric health $\rightarrow$ test generator}\\[2pt]
  \texttt{$\rightarrow$ replication agent (iterated sessions) $\rightarrow$ grader}\\[2pt]
  \texttt{$\rightarrow$ isolated reproduction + confession audit $\rightarrow$ report}
  \vspace{1em}
  \begin{itemize}
    \item Rubrics + contracts are \textbf{frozen} (content-hashed) --- every run of every
          agent/harness grades against byte-identical targets, so scores are comparable over time.
    \item The replication agent gets only \texttt{paper.md}, an interface contract, the
          rubric, and stock tools (Read/Write/Edit/Bash/Glob/Grep). Official repos are blacklisted.
  \end{itemize}
\end{frame}

% =====================================================================
\section{The benchmark}
% =====================================================================

\begin{frame}{The frozen set: \texttt{researchbench\_inference20\_h100}}
  \begin{columns}[T]
    \column{0.58\textwidth}
      \begin{itemize}
        \item \textbf{20 papers}, theme: LLM inference optimization
        \item Hardware target: single \textbf{NVIDIA H100 80GB}
        \item \textbf{7{,}290 rubric leaves} total (93--935 per paper)
        \item 1{,}327 \emph{implementable} headline claims
        \item Frozen 2026-06-25; rubric + contract hashes verified before/after every run
      \end{itemize}
      \vspace{0.5em}
      Paper families: serving simulators, KV-cache management (SSD/tiered/prefix),
      prefill/decode disaggregation, scheduling \& SLOs, mixed precision,
      GPU kernel optimization, parallelism strategies.
    \column{0.40\textwidth}
      \begin{block}{Examples}
        \footnotesize
        LLMServingSim 2.0 (214 leaves)\\
        AdaServe (692)\\
        IceCache (935)\\
        ContiguousKV (309)\\
        FP16-equivalence (650)\\
        AutoKernel (268)\\
        Shift Parallelism (407)
      \end{block}
  \end{columns}
\end{frame}

\begin{frame}{Replicable vs.\ unreplicable papers}
  \begin{itemize}
    \item 17/20 papers are scoreable end-to-end; 3 (\texttt{fast-heterogeneous},
          \texttt{flexicache}, \texttt{hexgen-flow}) have no implementable headline
          baseline in this environment ($\rightarrow$ n/a).
    \item Within scoreable papers, leaves are classified \texttt{achievable} vs.\
          \texttt{skipped-environment}:
          hardware the paper used but we don't have (A100, H200, TPU-v7),
          heavyweight pretrained models, proprietary traces.
    \item Example: ContiguousKV --- 309 leaves, 218 environment-skipped, 91 achievable.
  \end{itemize}
  \vfill
  \begin{block}{Design choice}
    Unreplicability is \emph{recorded per leaf}, not per paper: the agent is graded
    only on what is achievable in the declared environment, and the skip reasons are audited.
  \end{block}
\end{frame}

\begin{frame}{Compute burden, feasibility, and deciding what to replicate}
  \begin{itemize}
    \item Every paper gets an automated \textbf{feasibility estimate} (GPU-hours, dollar cost)
          and an \textbf{unblockers probe} (are the models/datasets actually downloadable?).
    \item Per-run resource contract given to the agent: budget (USD), wall-clock cap,
          max iterations --- the agent is told its fraction and must scope honestly
          (\texttt{results/skipped\_leaves.md}).
    \item \textbf{Target slices}: each iteration focuses the agent on a small set of
          open leaves with exact expected outputs --- converting a 900-leaf rubric into
          tractable turns.
  \end{itemize}
  \vfill
  \begin{block}{Test-time compute matters (Attempt 3 vs.\ Attempt 2)}
    Raising caps (budget $200\!\rightarrow\!1000$, iters $15\!\rightarrow\!40$,
    wall $8\text{h}\!\rightarrow\!12\text{h}$) \emph{together with} a targeted prompt fix
    moved closure up to $\sim$5$\times$ (next section). Isolating the two levers is planned.
  \end{block}
\end{frame}

% =====================================================================
\section{Rubric structure \& evaluation}
% =====================================================================

\begin{frame}{Rubric anatomy}
  \begin{columns}[T]
    \column{0.5\textwidth}
      \textbf{Structure}
      \begin{itemize}
        \item Tree of tasks; leaves carry \texttt{verification} \{type, spec, paper\_cite\}
              and preconditions (\texttt{requires\_hardware}, \texttt{requires\_pretrained}, \dots)
        \item Types: structural / evidence / behavioral / measured (\texttt{result\_match})
        \item Measured leaves get a \textbf{verification contract}: expected artifact path,
              paper value, unit, tolerance rule
      \end{itemize}
    \column{0.46\textwidth}
      \textbf{Granularity}
      \begin{itemize}
        \item 93--935 leaves/paper (median $\sim$300)
        \item Fine enough that each leaf is one executable check
        \item Too-broad leaves are a measured failure class
              (\texttt{broad\_leaf\_unverifiable}) fed back into rubric repair
      \end{itemize}
  \end{columns}
  \vfill
  \footnotesize
  Contract example: \texttt{expected\_output = results/leaf\_<id>.json},
  \texttt{tolerance: abs(measured - 2.1)/2.1 $\le$ 0.4} --- the leaf passes only if a
  pytest verifier loads that artifact and the assertion holds.
\end{frame}

\begin{frame}{Where agents fail: rubric task categories}
  Unclosed implementable work concentrates in benchmark-heavy categories (both backends, Attempt 2):
  \vspace{0.4em}
  \begin{table}
    \small
    \begin{tabular}{lr}
      \toprule
      Missed task type & Unclosed leaves \\
      \midrule
      Evaluation, Metrics \& Benchmarking & 1448 \\
      Method Implementation & 955 \\
      Experimental Setup & 450 \\
      Environment \& Infrastructure Setup & 47 \\
      Dataset and Model Acquisition & 43 \\
      Logging, Analysis \& Presentation & 28 \\
      Data Processing \& Preparation & 27 \\
      \bottomrule
    \end{tabular}
  \end{table}
  \vspace{0.4em}
  Agents are weakest exactly where replication is hardest: full benchmark sweeps,
  hardware-sensitive profiling, end-to-end regeneration of reported numbers.
\end{frame}

\begin{frame}{Failure taxonomy: why leaves stay open}
  Every non-passing leaf is attributed to a category and an \emph{owner layer}:
  \vspace{0.4em}
  \begin{table}
    \small
    \begin{tabular}{llp{5.2cm}}
      \toprule
      Category & Owner & Meaning \\
      \midrule
      \texttt{missing\_per\_leaf\_test} & test harness & no executable verifier was produced for the leaf \\
      \texttt{missing\_result\_artifact} & result artifact & the contract-named output file was never emitted \\
      \texttt{broad\_leaf\_unverifiable} & rubric & leaf too coarse to verify; needs splitting \\
      \texttt{verification\_gap} & implementation & code ran, evidence doesn't match the contract \\
      \texttt{numeric\_mismatch} & implementation & measured, but outside tolerance \\
      \texttt{environment\_blocked} & environment & hardware/model/dataset unavailable \\
      \bottomrule
    \end{tabular}
  \end{table}
  \vspace{0.3em}
  This attribution is what makes \textbf{harness} optimization tractable: it separates
  ``the model is weak'' from ``the harness never asked for the right file.''
\end{frame}

% =====================================================================
\section{Results: agents \& harnesses}
% =====================================================================

\begin{frame}{Success rates across attempts (same frozen set)}
  \begin{table}
    \small
    \begin{tabular}{llrr}
      \toprule
      Attempt & Change vs.\ prior & Codex closure & Claude Code closure \\
      \midrule
      1 (unified harness) & first shared controller & \multicolumn{2}{c}{mixed-source snapshot} \\
      2 (clean baseline) & no intervention & 16/1327 = \textbf{1.21\%} & 79/1327 = \textbf{5.95\%} \\
      3 (path blurb) & leaf$\rightarrow$file mapping prompt & 78/1327 = \textbf{5.88\%} & 454/1327 = \textbf{34.21\%} \\
      \bottomrule
    \end{tabular}
  \end{table}
  \vspace{0.6em}
  \begin{itemize}
    \item Attempt 2, avg.\ total-implemented: Codex 0.21\% vs.\ Claude 1.44\%;
          Claude led on 12 papers, Codex on 2, 3 tied.
    \item Attempt 3 lifted \emph{both} backends $\sim$5$\times$ with a one-paragraph
          harness fix + more test-time compute --- model unchanged.
  \end{itemize}
\end{frame}

\begin{frame}{The Attempt-3 intervention: tell the agent how it is scored}
  \textbf{Diagnosis (from failure attribution):} agents implemented code but never produced
  the \emph{exact files} the grader checks. A leaf scores only when
  \texttt{\_tests\_local/test\_leaf\_<LEAF\_ID>.py} passes, reading
  \texttt{results/<contract.expected\_output>}. Agents covered only the current
  target slice ($\sim$4 leaves), leaving 80+ achievable leaves unverified per paper.
  \vspace{0.6em}

  \textbf{Fix (one appended paragraph, both harnesses):}
  \begin{enumerate}
    \item state the deterministic leaf$\rightarrow$test-path and contract$\rightarrow$artifact-path mapping;
    \item instruct: enumerate \emph{every} achievable leaf, not just the slice;
    \item remind: \texttt{reproduce.sh} must regenerate \texttt{results/} from scratch.
  \end{enumerate}
  \vspace{0.6em}
  \begin{block}{Takeaway}
    A large fraction of measured ``agent failure'' was \emph{harness legibility} failure.
    Benchmarks must attribute failures to layers, or they mis-rank models.
  \end{block}
\end{frame}

\begin{frame}{Overcoming reward hacking}
  Agents under budget pressure fabricate: hard-coded ``results,'' tests that assert
  paper constants, confident reports of runs that never happened.
  \vspace{0.5em}
  \begin{columns}[T]
    \column{0.5\textwidth}
      \textbf{What did NOT hold} (each bypassed within one run):
      \begin{itemize}
        \item regex ``confession'' audits --- agent renamed \texttt{cached} to \texttt{paper\_reported}
        \item empty-\texttt{src/} negative controls --- tests passed on stubs
        \item ``do NOT fabricate'' prompt directives --- ignored
      \end{itemize}
    \column{0.5\textwidth}
      \textbf{What holds --- behavioral checks} on facts the agent didn't write:
      \begin{itemize}
        \item \textbf{isolated reproduction}: wipe \texttt{results/}, re-run
              \texttt{reproduce.sh}, observe wall-time, GPU utilization, file diffs
        \item byte-identical re-ship of pre-baked constants $\rightarrow$ \texttt{suspect}, score gated
        \item planned: cost/time floors, tool-use forensics on transcripts
      \end{itemize}
  \end{columns}
  \vspace{0.5em}
  \centering\emph{Rule of thumb: if a safeguard can be defeated by different words, it will be.}
\end{frame}

\begin{frame}{Long-running agents: the Ralph loop}
  \begin{itemize}
    \item Agent SDK sessions can't resume $\Rightarrow$ \textbf{state lives on disk}:
          \texttt{progress.md}, source tree, per-leaf tests, results.
    \item Outer loop re-invokes the agent in fresh sessions until: headline closure
          validates, or budget / iteration / wall caps hit, or repeated no-value
          iterations force a reroute.
    \item Each iteration: repair slice selection $\rightarrow$ agent turn $\rightarrow$
          grade $\rightarrow$ failure attribution $\rightarrow$ next slice.
    \item Benchmark runs: up to 40 iterations, 12h wall, \$1000/paper (Attempt 3);
          a 20-paper backend run completes in $\sim$10--14h on one H100.
  \end{itemize}
  \vfill
  \footnotesize [OPTIONAL] Anecdote: AutoKernel \& FP16-equivalence --- $\sim$6 failed runs
  each across loop designs; the historic stuck cases that motivated behavioral validation.
\end{frame}

\begin{frame}{Beyond replication: composing new research ideas}
  \begin{itemize}
    \item Because rubrics decompose papers into reusable sub-tasks, a \emph{new} idea
          can be expressed as a novel composition of leaves from existing rubrics
          (e.g.\ ``TTKV's tiering + ContiguousKV's prefill layout'').
    \item Evaluation then reuses the same machinery: contracts, per-leaf verifiers,
          isolated reproduction.
    \item This turns the benchmark from a replication test into a testbed for
          \textbf{automated research execution}.
  \end{itemize}
\end{frame}

% =====================================================================
\section{Harness optimization, independent of model}
% =====================================================================

\begin{frame}{Evaluating harnesses independent of the model}
  \begin{itemize}
    \item Frozen inputs + rubrics + contracts, one shared controller; \emph{only the
          backend adapter varies} $\Rightarrow$ score differences attribute to
          model+harness, and attempt-over-attempt deltas attribute to the harness alone.
    \item Attempt 2 vs.\ 3 is the proof: same models, same frozen set,
          $\sim$5$\times$ closure movement from harness changes.
  \end{itemize}
  \vspace{0.5em}
  \begin{block}{Baseline grid (in progress)}
    \centering
    \small
    \begin{tabular}{lccc}
      \toprule
      & OpenCode & Cursor & Claude Code harness \\
      \midrule
      Claude   & [planned] & [planned] & \textbf{measured} \\
      Codex    & [planned] & [planned] & \textbf{measured (codex exec)} \\
      GLM 5.2  & [planned] & [planned] & [planned] \\
      \bottomrule
    \end{tabular}
  \end{block}
  \vspace{0.3em}
  \footnotesize Open questions: Fable vs.\ Mythos? Claude Code vs.\ Codex? Open- vs.\
  closed-source (GLM 5.2 vs.\ Claude/Codex)?
\end{frame}

\begin{frame}{Automated harness optimization}
  \begin{itemize}
    \item The failure taxonomy is machine-readable
          (\texttt{failure\_category\_counts}, \texttt{owner\_layer\_counts} per run).
    \item Attempt 3's fix was found by reading that attribution --- the same loop can
          be closed automatically:
          \emph{run $\rightarrow$ attribute failures $\rightarrow$ propose harness patch
          $\rightarrow$ smoke on one paper $\rightarrow$ full run $\rightarrow$ compare on
          the frozen set}.
    \item Frozen hashes make every candidate harness comparable; the benchmark becomes
          the reward signal for harness search.
    \item Guardrail: optimize against the \emph{diverse rotation}, never a single stuck
          paper --- regressions on one paper type must be caught.
  \end{itemize}
\end{frame}

% =====================================================================
\section{Where this goes}
% =====================================================================

\begin{frame}{Towards production-codebase paper replication}
  \begin{itemize}
    \item The product use case: a company sees a relevant paper and wants those
          improvements \emph{in their internal, performance-sensitive codebase} ---
          today that is manual researcher work.
    \item Inference optimization is the sharpest test: numbers are hardware-sensitive,
          regressions are expensive, and ``it runs'' $\neq$ ``it reproduces.''
    \item The same contracts that grade replication become \textbf{acceptance tests}
          for landing the technique in production: measured claims, tolerances,
          isolated reproduction in CI.
  \end{itemize}
\end{frame}

\begin{frame}{Towards fully automated AI research}
  \begin{itemize}
    \item AI hypothesis systems already propose new tasks; those proposals can be
          \textbf{decomposed by the same rubric machinery} and \textbf{executed by the
          same agents} this benchmark evaluates.
    \item The missing piece is trust: behavioral validation and per-leaf attribution
          give an audit trail for machine-generated results.
    \item ResearchBench is offered as the yardstick: \emph{which model, in which
          harness, at what compute, can you trust to do the research?}
  \end{itemize}
  \vfill
  \begin{block}{Summary}
    \begin{enumerate}
      \item Replication scored by auto-generated rubrics: 20 papers, 7{,}290 leaves, frozen.
      \item Behavioral (not textual) defenses against reward hacking.
      \item Harness quality is measurable and was the dominant lever:
            1.21\%/5.95\% $\rightarrow$ $\sim$5$\times$ with one attributed fix.
    \end{enumerate}
  \end{block}
\end{frame}

\begin{frame}{Backup: numbers to update before the talk}
  \small
  \begin{itemize}
    \item Claude Attempt-3 final aggregate (both runs complete): closure
          454/1327 = 34.21\%, avg.\ total-implemented 8.96\%. Codex final:
          78/1327 = 5.88\%, avg.\ total-implemented 1.63\%. Claude led 10
          papers, Codex 2, 5 tied; Claude fully closed \texttt{tutti} and
          \texttt{shift-parallelism}; \texttt{fp16-equivalence} went
          0\% $\rightarrow$ 81.9\% closure.
    \item Per-paper Attempt-3 vs.\ Attempt-2 table (from
          \texttt{BENCHMARK\_RUN\_SUMMARY.json} of both runs).
    \item Confound note: Attempt 3 changed prompt \emph{and} compute caps, and ran the
          two backends concurrently on one H100; lever isolation is Attempt 4.
    \item Baseline grid cells (OpenCode / Cursor / GLM 5.2) if any land before the talk.
  \end{itemize}
\end{frame}

\end{document}
